% =====================================================================
%  第 20 章  重夸克物理：B 介子衰变与重夸克偶素            （ch20_heavy.tex，writer W23）
%  约定：本书约定（WRITER_SPEC_v2 §3）；n_A=(1^+,0^-,0_T)、n_B=(0^+,1^-,0_T)（不用 SCET 的 n、nbar）；
%  D_mu = d_mu - i g A_mu^a t^a；D = 4-2 eps，S_eps=(4pi)^eps e^{-eps gamma_E}；
%  BBL、Collins 用 D = d + i g A（g -> -g）。验证脚本：writing/work/W23/s1_bdpi.py … s7_extra.py。
% =====================================================================
\chapter{重夸克物理：$B$ 介子衰变与重夸克偶素}\label{ch:heavy}

重夸克质量 $m_Q\gg\Lqcd$ 本身就是硬标度：$b$ 夸克衰变使末态轻强子以接近光速运动，出现与硬遍举过程（\cref{ch:excl}）相同的共线结构；重夸克偶素中相对速度 $v\ll1$ 提供另一个展开参数。衰变率是固定物理点上的 Minkowski 量，把 OPE 或 NRQCD 衰变公式用于固定运动学点上的强子衰变率还需要局域夸克—强子对偶（\cref{conj:ee-duality}，[C]），这不在任何微扰证明的范围之内（见下框）。

\begin{tcolorbox}[enhanced,breakable,title=本章证明了什么／没有证明什么,fonttitle=\bfseries,
  colback=gray!3!white,colframe=gray!55!black,boxrule=0.6pt,arc=2pt,
  left=5pt,right=5pt,top=3pt,bottom=3pt,before skip=6pt,after skip=6pt]
\textbf{证明了：}\cref{prop:heavy-nrqcd-matching}（\tagR，$R_{\mathrm{PT}}$：电磁衰变匹配系数的树图值与单圈 $c_v$）。\par
\textbf{给出全阶论证：}\cref{thm:heavy-bdpi,prop:heavy-btopi-ff}（\tagA；瓶颈：\cref{rem:scet-assumptions}的 SCET 前提、端点压低假设）；\cref{thm:heavy-inclusive-b}（\tagA；瓶颈：\cref{thm:dis-ope}与 SCET 前提）；\cref{thm:heavy-nrqcd-decay,prop:heavy-quarkonium-fragmentation}（\tagA）；强子层次均另需\cref{asm:intro-bridge}。\par
\textbf{没有证明：}无粲两体衰变领头幂次因子化的全阶形式（\cref{prop:heavy-qcdf-lp}，\tagP）；其幂次修正的标准共线因子化（\cref{prop:heavy-qcdf-endpoint}，\tagV：卷积发散）；产生的 NRQCD 因子化（\cref{conj:heavy-nrqcd-production}，\tagC）；固定运动学点上所需的局域对偶（\cref{conj:ee-duality}，\tagC）。
\end{tcolorbox}

\paragraph{记号。}
用 $n_A=(1^+,0^-,\mathbf 0_T)$、$n_B=(0^+,1^-,\mathbf 0_T)$ 而不用 SCET 的 $n$、$\bar n$（引用\cref{ch:scet}的文献时 $n\to\sqrt2\,n_A$、$\bar n\to\sqrt2\,n_B$）；被发射的轻介子或强子喷注总沿 $+z$ 运动，$p\simeq p^+n_A$，其共线场经 $n_A\cdot A_s$ 与软胶子耦合，共线 Wilson 线沿 $n_B$。$\lambda=\Lqcd/Q$；HQET 场 $h_v(x)=\ee^{\ii m_Qv\cdot x}\frac{1+\sla v}{2}Q(x)$；BBL 与 Collins 的公式已按 $g\to-g$ 换算（\cref{def:prelim-qcd}）。
